% Mathématiques 3e — V3 LaTeX
% Triangles semblables et démonstrations

\section{Objectifs}
\begin{itemize}
\item Reconnaître et caractériser des triangles semblables.
\item Utiliser des rapports de côtés correspondants.
\item Mobiliser les angles pour justifier une similitude.
\item Construire une démonstration géométrique structurée.
\end{itemize}

\section{Repères et vocabulaire}
Deux triangles semblables ont la même forme mais pas nécessairement la même taille. En troisième, la similitude relie proportionnalité, angles, Thalès et transformations. Elle fournit aussi un terrain privilégié pour apprendre à organiser une démonstration.

\section{1. Définition par les côtés}
Deux triangles sont semblables lorsque leurs côtés correspondants sont proportionnels.

\[
\dfrac{AB}{A'B'}=\dfrac{BC}{B'C'}=\dfrac{CA}{C'A'}
\]

\begin{exemple}
L’ordre des sommets est essentiel : il fixe quelles longueurs et quels angles se correspondent.
\end{exemple}

\section{2. Caractérisation angulaire}
Deux triangles ayant leurs angles correspondants égaux ont la même forme et sont semblables.

\[
\widehat A=\widehat{A'}
\]

\[
\widehat B=\widehat{B'}
\]

\begin{exemple}
Comme la somme des angles d’un triangle vaut $180^\circ$, deux égalités d’angles suffisent pour obtenir la troisième.
\end{exemple}

\coursefigure{/images/mathematiques/3eme/triangles-semblables-demonstrations-1.svg}{Première représentation structurante pour Triangles semblables et démonstrations.}{La figure sert de support visuel pour organiser les relations géométriques ou algorithmiques du chapitre.}

\section{3. Coefficient de similitude}
Le rapport commun entre longueurs correspondantes joue le rôle d’un coefficient d’agrandissement ou de réduction.

\[
k=\dfrac{A'B'}{AB}
\]

\begin{exemple}
Si $k>1$, le second triangle est plus grand ; si $0<k<1$, il est plus petit.
\end{exemple}

\section{4. Aires de triangles semblables}
Lorsque les longueurs sont multipliées par $k$, les aires sont multipliées par $k^2$.

\[
A'=k^2A
\]

\begin{exemple}
Cette relation permet de retrouver un coefficient de similitude à partir d’aires.
\end{exemple}

\section{5. Cas d’égalité}
Des triangles égaux sont un cas particulier de triangles semblables avec coefficient $k=1$.

\[
k=1
\]

\begin{exemple}
Les cas d’égalité permettent de prouver des longueurs ou angles égaux dans des configurations où aucune mesure numérique n’est donnée.
\end{exemple}

\coursefigure{/images/mathematiques/3eme/triangles-semblables-demonstrations-2.svg}{Deuxième représentation pédagogique pour Triangles semblables et démonstrations.}{La seconde figure permet de comparer des configurations, des états ou des transformations dans les problèmes du chapitre.}

\section{6. Parallélisme et angles}
Des droites parallèles créent des égalités d’angles correspondants ou alternes-internes, ce qui peut établir la similitude de deux triangles.

\[
(d_1)\parallel(d_2)
\]

\begin{exemple}
On doit citer la propriété d’angles utilisée et identifier les droites concernées.
\end{exemple}

\section{7. Thalès et similitude}
Les configurations de Thalès produisent des triangles semblables : les rapports de longueurs sont une conséquence de cette même structure géométrique.

\[
\dfrac{AM}{AB}=\dfrac{AN}{AC}=\dfrac{MN}{BC}
\]

\begin{exemple}
La similitude donne un point de vue global sur le théorème sans remplacer la vérification de ses hypothèses.
\end{exemple}

\section{8. Rédiger une preuve}
Une démonstration distingue données, propriété utilisée et conclusion. Le dessin aide à chercher mais ne remplace pas l’argument.

\[
\text{données}\Rightarrow\text{propriété}\Rightarrow\text{conclusion}
\]

\begin{exemple}
Une rédaction efficace est courte mais explicite : elle cite les faits nécessaires et évite les affirmations « parce que cela se voit ».
\end{exemple}

\section{Méthode générale de résolution}
\begin{methode}
\begin{enumerate}
\item Identifier les sommets correspondants.
\item Relever les égalités d’angles ou les rapports de longueurs disponibles.
\item Choisir la caractérisation la plus directe.
\item Écrire les rapports dans un ordre cohérent.
\item Effectuer le calcul ou la déduction géométrique.
\item Conclure précisément sur la similitude, une longueur ou un angle.
\end{enumerate}
\end{methode}

\section{Problème guidé et stratégie}
Un exercice de troisième peut demander plusieurs décisions successives : reconnaître une configuration, sélectionner une propriété, produire une valeur exacte ou approchée, puis interpréter. On commence par écrire les hypothèses et la grandeur recherchée. On ne remplace par des nombres qu’après avoir écrit la relation mathématique ou logique adaptée. La conclusion reprend le vocabulaire de l’énoncé et précise l’unité ou la propriété démontrée.

\begin{attention}
Une construction visuelle ou un résultat de calculatrice peut suggérer une réponse, mais une preuve géométrique, un raisonnement algébrique ou une analyse d’algorithme doit expliciter pourquoi la conclusion est valide.
\end{attention}

\section{Contrôler un résultat}
Le contrôle dépend du chapitre : comparer les longueurs dans une figure, vérifier qu’un angle reste dans l’intervalle attendu, tester une valeur dans une équation, suivre l’état d’un programme ou convertir l’unité finale. Une deuxième méthode ou un cas simple permet souvent de détecter une erreur avant la réponse définitive.

\section{Erreurs fréquentes}
\begin{itemize}
\item Comparer des côtés non correspondants.
\item Déduire la similitude d’un dessin à l’échelle.
\item Inverser un seul rapport et pas les autres.
\item Confondre triangles égaux et triangles semblables.
\item Utiliser une égalité d’angles sans justifier le parallélisme ou la propriété invoquée.
\item Donner un coefficient d’aire égal au coefficient de longueur.
\end{itemize}

\section{À retenir}
\begin{aretenir}
\begin{itemize}
\item Les triangles semblables ont des côtés correspondants proportionnels.
\item Ils ont des angles correspondants égaux.
\item Le coefficient $k$ agit sur les longueurs et $k^2$ sur les aires.
\item La similitude relie proportionnalité, Thalès et transformations.
\end{itemize}
\end{aretenir}

\section{Réinvestissement : Coefficient de similitude}
Cette notion peut être réinvestie dans un problème où plusieurs informations doivent être reliées. Le rapport commun entre longueurs correspondantes joue le rôle d’un coefficient d’agrandissement ou de réduction. L’élève commence par expliciter les données pertinentes, puis choisit une propriété et contrôle sa conclusion. Si $k>1$, le second triangle est plus grand ; si $0<k<1$, il est plus petit. Cette étape de réinvestissement prépare les problèmes de brevet où la stratégie compte autant que le calcul.

\section{Réinvestissement : Définition par les côtés}
Cette notion peut être réinvestie dans un problème où plusieurs informations doivent être reliées. Deux triangles sont semblables lorsque leurs côtés correspondants sont proportionnels. L’élève commence par expliciter les données pertinentes, puis choisit une propriété et contrôle sa conclusion. L’ordre des sommets est essentiel : il fixe quelles longueurs et quels angles se correspondent. Cette étape de réinvestissement prépare les problèmes de brevet où la stratégie compte autant que le calcul.

\section{Réinvestissement : Parallélisme et angles}
Cette notion peut être réinvestie dans un problème où plusieurs informations doivent être reliées. Des droites parallèles créent des égalités d’angles correspondants ou alternes-internes, ce qui peut établir la similitude de deux triangles. L’élève commence par expliciter les données pertinentes, puis choisit une propriété et contrôle sa conclusion. On doit citer la propriété d’angles utilisée et identifier les droites concernées. Cette étape de réinvestissement prépare les problèmes de brevet où la stratégie compte autant que le calcul.
